\question
Дверь с двумя замками. Моргана нашёл запертую дверь. Футаба подключила к панели восемь независимых двоичных сигналов $x_1,\ldots,x_8$; единица означает активный сигнал. Дверь открывается ровно при выполнении всех условий:
\begin{enumerate}
\item первый замок открывается, если активны оба сигнала хотя бы одной из пар $(x_1,x_2)$ и $(x_3,x_4)$;
\item второй замок при $x_5=0$ требует разных значений $x_6,x_7$, а при $x_5=1$ --- $x_7=x_8=1$; остальные сигналы на этот замок не влияют.
\end{enumerate}
Составьте булеву функцию открытия двери и схему из двухвходовых элементов $\land,\lor,\oplus$ и одновходовых $\neg$. Разветвление проводов разрешено; константных входов нет. Условия полностью задают поведение панели на всех наборах сигналов.

---------------

Автор -- Олег Эмрих
